% ======================================================================
\chapter[Кварковая структура адронов. Квантовая хромодинамика]%
        {Кварковая структура адронов. Сильное взаимодействие и КХД}
\label{t:15}
% ======================================================================

\begin{qbox}
\textbf{Билет 15.} Кварковая структура адронов — мезоны, барионы. Сильное взаимодействие,
квантовая хромодинамика.
\end{qbox}

\section{Содержание билета}

\textbf{Подтемы:}
\begin{enumerate}
 \item адроны; кварковая модель: мезоны ($q\bar q$) и барионы ($qqq$); квантовые числа и их проверка;
 \item цвет: проблема $\Delta^{++}$ и $\Omega^-$, бесцветность адронов, экспериментальная проверка числа цветов (отношение $R$);
 \item экспериментальные свидетельства кварков: глубоконеупругое рассеяние, струи;
 \item квантовая хромодинамика: глюоны, самодействие, асимптотическая свобода, конфайнмент, потенциал кварк--антикварк;
 \item сильное взаимодействие нуклонов как «остаточное» цветовое; происхождение масс адронов.
\end{enumerate}

\textbf{Необходимые понятия:} Стандартная модель, кварки и переносчики (\tref{t:14}); принцип Паули
(\tref{t:13}); законы сохранения зарядов; спин и сложение моментов.

\section{Определения и физические принципы}

\begin{definition}[адроны]
\emph{Адроны} — частицы, участвующие в сильном взаимодействии. Делятся на \emph{мезоны} (целый спин,
барионное число $B=0$) и \emph{барионы} (полуцелый спин, $B=1$; антибарионы — $B=-1$). Лёгчайшие барионы —
нуклоны (протон и нейтрон), лёгчайшие мезоны — $\pi$-мезоны.
\end{definition}

\begin{principle}{Кварковая модель (Гелл-Ман, Цвейг, 1964)}
Мезоны — связанные состояния кварка и антикварка ($q\bar q'$), барионы — трёх кварков ($qqq$).
Кварки имеют спин $\frac12$, барионное число $\frac13$, электрический заряд $+\frac23e$ ($u,c,t$) или $-\frac13e$ ($d,s,b$).
Квантовые числа адрона — суммы квантовых чисел составляющих (с учётом орбитального движения для спина и чётности).
\end{principle}

\begin{definition}[цвет и квантовая хромодинамика]
\emph{Цвет} — квантовое число кварков, принимающее три значения ($r,g,b$); антикварки несут антицвета.
\emph{Квантовая хромодинамика} (КХД) — калибровочная теория сильного взаимодействия с группой $SU(3)$:
кварки взаимодействуют обменом восемью безмассовыми глюонами, которые сами несут цвет.
Наблюдаемые адроны — цветовые синглеты («бесцветны»).
\end{definition}

\emph{Физический смысл.} Сильное взаимодействие нуклонов в ядре (\tref{t:17}) — не фундаментальное, а
«остаточное» проявление цветового взаимодействия кварков, подобно ван-дер-ваальсовым силам между
нейтральными атомами как остатку электромагнитных сил.

\emph{Область применимости.} Кварковая модель с «составляющими» кварками качественно описывает спектр и
квантовые числа адронов. Теория возмущений КХД применима лишь при больших передачах импульса
($Q\gg1$~ГэВ), где $\alpha_s\ll1$; явления на масштабе $\sim1$~фм (массы адронов, конфайнмент) требуют
непертурбативных методов (решёточная КХД).

\section{Основная теория}

\subsection{Проблема: систематика адронов}

К началу 1960-х годов было известно множество адронов. Группировка частиц с одинаковыми спином и чётностью
в координатах «проекция изоспина -- странность» дала правильные фигуры — октеты и декуплет («восьмеричный
путь», Гелл-Ман, Нееман, 1961). Такие мультиплеты естественно возникают, если адроны составлены из трёх
лёгких кварков $u,d,s$ (\emph{странность} $S=-1$ для $s$-кварка, $+1$ для $\bar s$). Предсказание по декуплету
бариона $\Omega^-$ (масса $\approx1680$~МэВ, $S=-3$) и его открытие в 1964~г. (масса 1672~МэВ) стало триумфом
модели.

\subsection{Мезоны}

Мезон $q\bar q'$: $B=\frac13-\frac13=0$; спины кварков $\frac12\otimes\frac12$ дают $S=0$ или $1$; при нулевом
орбитальном моменте — псевдоскалярные ($J=0$) и векторные ($J=1$) мезоны.
\begin{center}
\small
\begin{tabular}{@{}lllll@{}}
\toprule
Мезон & Состав & $Q/e$ & $J$ & Масса, МэВ \\
\midrule
$\pi^+$, $\pi^-$ & $u\bar d$, $d\bar u$ & $+1$, $-1$ & 0 & 139,6 \\
$\pi^0$ & $(u\bar u-d\bar d)/\sqrt2$ & 0 & 0 & 135,0 \\
$K^+$, $K^0$ & $u\bar s$, $d\bar s$ & $+1$, 0 & 0 & 493,7; 497,6 \\
$\rho^+$ & $u\bar d$ & $+1$ & 1 & 775 \\
$J/\psi$ & $c\bar c$ & 0 & 1 & 3097 \\
$\Upsilon$ & $b\bar b$ & 0 & 1 & 9460 \\
\bottomrule
\end{tabular}
\end{center}
Открытие узкого резонанса $J/\psi$ в 1974~г. («ноябрьская революция») доказало существование $c$-кварка;
$\Upsilon$ (1977) — $b$-кварка; $t$-кварк открыт в 1995~г. на Тэватроне (он распадается раньше, чем успевает
образовать адрон).

\subsection{Барионы}

Барион $qqq$: $B=1$; спины трёх кварков дают $J=\frac12$ или $\frac32$ (при нулевых орбитальных моментах).
\begin{center}
\small
\begin{tabular}{@{}lllll@{}}
\toprule
Барион & Состав & $Q/e$ & $J$ & Масса, МэВ \\
\midrule
$p$ & $uud$ & $+1$ & $\frac12$ & 938,3 \\
$n$ & $udd$ & 0 & $\frac12$ & 939,6 \\
$\Lambda^0$ & $uds$ & 0 & $\frac12$ & 1116 \\
$\Sigma^+$ & $uus$ & $+1$ & $\frac12$ & 1189 \\
$\Delta^{++}$ & $uuu$ & $+2$ & $\frac32$ & 1232 \\
$\Omega^-$ & $sss$ & $-1$ & $\frac32$ & 1672 \\
\bottomrule
\end{tabular}
\end{center}
Проверка зарядов (вывод~\ref{der:15:charges}): $p$: $\frac23+\frac23-\frac13=1$; $n$: $\frac23-\frac13-\frac13=0$;
$\Delta^{++}$: $3\cdot\frac23=2$; $\Omega^-$: $3\cdot(-\frac13)=-1$. Нет адронов с дробным зарядом и нет
адронов из двух кварков — модель объясняет, почему.

\begin{figure}[htbp]
\centering
\begin{tikzpicture}[scale=1.05]
  \foreach \x/\lab in {0/{$\Delta^-$},1/{$\Delta^0$},2/{$\Delta^+$},3/{$\Delta^{++}$}} { \fill[accent] (\x,0) circle (3pt); \node[above, font=\small] at (\x,0.1) {\lab}; }
  \foreach \x/\lab in {0.5/{$\Sigma^{*-}$},1.5/{$\Sigma^{*0}$},2.5/{$\Sigma^{*+}$}} { \fill[accent] (\x,-1) circle (3pt); \node[above, font=\small] at (\x,-0.9) {\lab}; }
  \foreach \x/\lab in {1/{$\Xi^{*-}$},2/{$\Xi^{*0}$}} { \fill[accent] (\x,-2) circle (3pt); \node[above, font=\small] at (\x,-1.9) {\lab}; }
  \fill[keycol] (1.5,-3) circle (3.5pt); \node[below, font=\small, keycol] at (1.5,-3.1) {$\Omega^-$ ($sss$)};
  \draw[gray] (0,0) -- (3,0) -- (1.5,-3) -- cycle;
  \foreach \y/\S/\m in {0/0/1232,-1/{-1}/1385,-2/{-2}/1533,-3/{-3}/1672} {
     \node[font=\small, anchor=west] at (4.0,\y) {$S=\S$, \ $M\approx\m$~МэВ};
  }
  \node[font=\scriptsize, below] at (0,-0.12) {$ddd$};
  \node[font=\scriptsize, below] at (3,-0.12) {$uuu$};
\end{tikzpicture}
\caption{Барионный декуплет $J=\frac32$. Каждый шаг вниз — замена $u$ или $d$ кварка на $s$ и увеличение массы
примерно на 150~МэВ; эквидистантность позволила предсказать массу $\Omega^-$.}
\label{fig:15:decuplet}
\end{figure}

\subsection{Цвет}

\paragraph{Проблема.} $\Delta^{++}=uuu$ со спином $\frac32$ и $L=0$: пространственная функция симметрична (основное
состояние), спиновая — тоже (все три спина параллельны), ароматная — тоже (три $u$). Полная функция трёх
тождественных фермионов оказывается симметричной — противоречие с принципом Паули. То же для $\Omega^-=sss$.

\paragraph{Решение} (Боголюбов, Струминский, Тавхелидзе; Хан, Намбу, 1965): кварки несут дополнительное
квантовое число с тремя значениями — цвет; волновая функция бариона антисимметрична по цвету:
\[
  \Psi_{\text{цвет}}=\frac{1}{\sqrt6}\sum_{ijk}\varepsilon_{ijk}\,q_iq_jq_k ,
\]
и полная функция (пространство $\times$ спин $\times$ аромат $\times$ цвет) становится антисимметричной. Мезон —
$\frac{1}{\sqrt3}(r\bar r+g\bar g+b\bar b)$. Обе комбинации инвариантны относительно цветовых преобразований
(цветовые синглеты). \emph{Гипотеза конфайнмента}: в природе наблюдаются только бесцветные состояния. Это
объясняет, почему существуют $q\bar q$ и $qqq$ (а также экзотические тетра- и пентакварки, обнаруженные в 2000--2010-х годах), но
не $qq$ или одиночные кварки.

\paragraph{Экспериментальная проверка числа цветов.} Отношение сечений (вывод~\ref{der:15:R})
\begin{equation}\label{eq:15:R}
  R=\frac{\sigma(e^+e^-\to\text{адроны})}{\sigma(e^+e^-\to\mu^+\mu^-)}=N_c\sum_qQ_q^2
\end{equation}
при энергиях, где рождаются $u,d,s$, равно $N_c\cdot\frac23$, выше порога $c\bar c$ — $N_c\cdot\frac{10}{9}$, выше $b\bar b$ —
$N_c\cdot\frac{11}{9}$. Измерения дают $R\approx2$, $\frac{10}{3}$, $\frac{11}{3}$ (с небольшой поправкой КХД $1+\alpha_s/\pi$), т.\,е.
$N_c=3$. Число цветов подтверждается также шириной распада $\pi^0\to2\gamma$ (амплитуда $\propto N_c$).

\subsection{Экспериментальные свидетельства кварков}

\begin{enumerate}
 \item \textbf{Глубоконеупругое рассеяние электронов на протонах} (SLAC, 1968; Нобелевская премия 1990): при больших
       передачах импульса сечение ведёт себя как при рассеянии на \emph{точечных} частицах со спином $\frac12$
       («партонах») — аналог опыта Резерфорда для протона. Выяснилось также, что кварки несут лишь около
       половины импульса протона; остальное — глюоны.
 \item \textbf{Адронные струи.} В $e^+e^-$-аннигиляции рождённые кварк и антикварк превращаются в два узких пучка адронов
       (струи), угловое распределение которых соответствует спину $\frac12$. Трёхструйные события (DESY, 1979) —
       испускание жёсткого глюона: прямое наблюдение глюона.
 \item \textbf{Спектроскопия} кваркония ($c\bar c$, $b\bar b$) — система уровней, подобная позитронию, хорошо описываемая
       потенциалом~\eqref{eq:15:V}.
\end{enumerate}

\subsection{Квантовая хромодинамика}

\paragraph{Глюоны.} Кварк при испускании глюона может менять цвет ($r\to g$), поэтому глюон несёт пару «цвет--антицвет»;
из $3\times3=9$ комбинаций одна (синглет $(r\bar r+g\bar g+b\bar b)/\sqrt3$) не меняет цвета и не входит в теорию, остаётся
$8$ глюонов. Глюоны цветные и взаимодействуют друг с другом (вершины $ggg$ и $gggg$) — в этом главное отличие
от фотона, который не заряжен.

\paragraph{Асимптотическая свобода.} Эффективная константа связи зависит от передаваемого импульса $Q$
(«бегущая константа»). В КЭД заряд экранируется виртуальными парами, и $\alpha$ медленно растёт с энергией. В КХД
самодействие глюонов приводит к \emph{антиэкранировке}, и $\alpha_s$ убывает (Гросс, Вильчек, Политцер, 1973;
Нобелевская премия 2004). В однопетлевом приближении (без вывода)
\begin{equation}\label{eq:15:alphas}
  \alpha_s(Q^2)=\frac{12\pi}{(33-2n_f)\ln(Q^2/\Lambda^2)},
\end{equation}
где $n_f$ — число активных ароматов, $\Lambda\approx0{,}2$--$0{,}3$~ГэВ. Значения: $\alpha_s(M_Zc)\approx0{,}118$, при $Q\sim1$~ГэВ
$\alpha_s\sim0{,}5$. На малых расстояниях кварки почти свободны (что и видно в глубоконеупругом рассеянии), на расстояниях
$\sim\hbar/\Lambda\sim1$~фм взаимодействие становится сильным.

\paragraph{Конфайнмент.} Феноменологический потенциал кварк--антикварк:
\begin{equation}\label{eq:15:V}
  V(r)=-\frac43\frac{\alpha_s\hbar c}{r}+\kappa r,\qquad \kappa\approx0{,}9\text{--}1\ \text{ГэВ/фм}.
\end{equation}
На малых $r$ — кулоноподобное притяжение (обмен одним глюоном; множитель $\frac43$ — цветовой), на больших —
линейно растущий потенциал: силовые линии глюонного поля собираются в узкую «трубку» (струну) с постоянной
энергией на единицу длины, и сила $\approx\kappa\approx16$~т не убывает с расстоянием. При растяжении струны на
$\sim1$--$1{,}5$~фм её энергия становится достаточной для рождения пары $q\bar q$, струна рвётся, и образуются два
бесцветных адрона — изолированный кварк не появляется (рис.~\ref{fig:15:V}). Этот процесс (адронизация) порождает струи.

\begin{figure}[htbp]
\centering
\begin{tikzpicture}
\begin{axis}[width=0.66\linewidth, height=5.0cm, axis lines=middle, xlabel={$r$, фм}, ylabel={$V$, ГэВ},
   xmin=0, xmax=2.1, ymin=-1.6, ymax=1.8, samples=200,
   xlabel style={at={(ticklabel* cs:1)}, anchor=west}, ylabel style={at={(ticklabel* cs:1)}, anchor=south}]
  \addplot[very thick, accent, domain=0.08:2.0] {-4/3*0.3*0.1973/x + 0.9*x - 0.4};
  \addplot[gray, dashed, domain=0.08:2.0] {-4/3*0.3*0.1973/x - 0.4};
  \addplot[gray, dotted, domain=0:2.0] {0.9*x-0.4};
  \node[font=\small, gray, anchor=west] at (axis cs:1.15,-0.85) {кулоновская часть};
  \node[font=\small, gray, anchor=west] at (axis cs:1.1,0.25) {$\kappa r$};
\end{axis}
\end{tikzpicture}
\caption{Потенциал кварк--антикварк~\eqref{eq:15:V} при $\alpha_s=0{,}3$, $\kappa=0{,}9$~ГэВ/фм (начало отсчёта энергии
выбрано произвольно). На больших расстояниях энергия растёт линейно — конфайнмент.}
\label{fig:15:V}
\end{figure}

\paragraph{Происхождение масс адронов.} Суммарная масса токовых кварков протона $2m_u+m_d\approx9$~МэВ — около 1\,\% его массы.
Остальное — энергия глюонного поля и кинетическая энергия кварков, локализованных в области $\sim1$~фм
($p\sim\hbar/R\sim200$--$300$~МэВ$/c$). В модели составляющих кварков им приписывают «эффективные» массы $\sim300$~МэВ
($u,d$) и $\sim500$~МэВ ($s$); это объясняет, например, $M_\Lambda-M_p\approx180$~МэВ.

\paragraph{Ядерные силы.} Нуклоны бесцветны, поэтому между ними действуют лишь остаточные силы на расстояниях, где
их кварковые облака перекрываются. На расстояниях $\sim1$--$2$~фм притяжение эффективно описывается обменом пионами
(лёгчайшими $q\bar q$ состояниями), радиус $\hbar/m_\pi c\approx1{,}4$~фм (\tref{t:14}).

\section{Математические выводы}

\begin{derivation}[label=der:15:charges]{квантовые числа адронов из кварков}
Используем аддитивность: $Q=\sum Q_q$, $B=\sum B_q$ ($B_q=\frac13$, $B_{\bar q}=-\frac13$), $S=-(N_s-N_{\bar s})$.

$\pi^+=u\bar d$: $Q=\frac23+\frac13=1$, $B=0$, $S=0$. \quad $K^+=u\bar s$: $Q=\frac23+\frac13=1$, $S=+1$.

$\Lambda=uds$: $Q=\frac23-\frac13-\frac13=0$, $B=1$, $S=-1$. \quad $\Omega^-=sss$: $Q=-1$, $B=1$, $S=-3$.

\emph{Спины.} Сложение трёх спинов $\frac12$: $\frac12\otimes\frac12\otimes\frac12=\frac32\oplus\frac12\oplus\frac12$; при $L=0$ барионы имеют
$J=\frac12$ (октет) или $\frac32$ (декуплет). Два спина: $\frac12\otimes\frac12=1\oplus0$ — мезоны $J=0$ и $1$.

\emph{Почему $\Delta^{++}$ требует цвета:} спин $\frac32$ означает $m_s=+\frac12$ у всех трёх кварков (в состоянии $J_z=\frac32$),
ароматы одинаковы, $L=0$ — функция без цвета симметрична по всем перестановкам, что запрещено для фермионов.
\end{derivation}

\begin{derivation}[label=der:15:R]{отношение $R$ и число цветов}
\emph{Допущения:} энергия $\sqrt s$ много больше масс рождающихся кварков, но ниже порога $Z^0$; аннигиляция идёт через
виртуальный фотон; поправки КХД отброшены.

Сечение $e^+e^-\to\gamma^*\to f\bar f$ для точечного фермиона с зарядом $Q_fe$ пропорционально квадрату заряда
(амплитуда $\propto eQ_fe$): $\sigma(f\bar f)=Q_f^2\,\sigma(\mu^+\mu^-)$, $\sigma(\mu^+\mu^-)=\dfrac{4\pi\alpha^2(\hbar c)^2}{3s}$.
Кварк каждого аромата рождается в $N_c$ цветовых вариантах, а адронизация происходит с вероятностью 1:
\[
  R=\frac{\sum_q N_c\,Q_q^2\,\sigma(\mu\mu)}{\sigma(\mu\mu)}=N_c\sum_qQ_q^2 .
\]
Для $u,d,s$: $\sum Q^2=\frac49+\frac19+\frac19=\frac23$; с $c$: $+\frac49=\frac{10}9$; с $b$: $+\frac19=\frac{11}9$. При $N_c=3$:
$R=2$; $\frac{10}{3}\approx3{,}3$; $\frac{11}{3}\approx3{,}7$, что согласуется с опытом. При $N_c=1$ было бы $R=\frac23$, $\frac{10}9$, $\frac{11}9$ —
втрое меньше наблюдаемого.
\end{derivation}

\begin{derivation}[label=der:15:string]{оценка расстояния разрыва струны}
Энергия струны длины $r$ равна $\kappa r$. Рождение лёгкой пары $q\bar q$ (с образованием двух мезонов) энергетически
выгодно, когда $\kappa r\gtrsim2m_{\text{эфф}}c^2$, где $m_{\text{эфф}}c^2\sim0{,}3$--$0{,}5$~ГэВ (масса составляющего кварка или
порядок массы лёгкого адрона). Отсюда $r\gtrsim2\cdot(0{,}3\text{--}0{,}5)/0{,}9$~фм $\approx0{,}7$--$1{,}1$~фм — масштаб размеров адронов.
Сила натяжения: $\kappa=1\ \text{ГэВ/фм}=1{,}6\cdot10^{-10}\ \text{Дж}/10^{-15}\ \text{м}=1{,}6\cdot10^{5}$~Н $\approx16$~тонн-силы.
\end{derivation}

\begin{derivation}[label=der:15:af]{масштаб $\Lambda$ и размер адронов}
По~\eqref{eq:15:alphas} $\alpha_s\to\infty$ при $Q\to\Lambda$: теория возмущений перестаёт работать на масштабе импульсов
$\sim\Lambda\approx0{,}2$--$0{,}3$~ГэВ$/c$, т.\,е. на расстояниях $r\sim\hbar/\Lambda=\hbar c/(\Lambda c)\approx197/250\ \text{фм}\approx0{,}8$~фм. Это
естественный размер адронов: КХД содержит размерный масштаб, возникающий из безразмерной константы
благодаря её «бегу» (размерная трансмутация).
\end{derivation}

\section{Итоговые формулы}

\keyformula{\text{мезон}=q\bar q',\quad \text{барион}=qqq,\quad Q_u=+\tfrac23,\ Q_d=Q_s=-\tfrac13,\quad B_q=\tfrac13}
 {$Q$ — заряд в единицах $e$, $B$ — барионное число.}
 {кварковая модель; квантовые числа аддитивны.}
 {состав и квантовые числа адронов; проверка законов сохранения в реакциях.}

\keyformula{R=\frac{\sigma(e^+e^-\to\text{адроны})}{\sigma(e^+e^-\to\mu^+\mu^-)}=N_c\sum_qQ_q^2}
 {$N_c$ — число цветов, сумма по ароматам, рождение которых энергетически доступно.}
 {энергии много выше порогов кварков и ниже $M_Z$; низший порядок по $\alpha_s$.}
 {экспериментальное определение $N_c=3$ и зарядов кварков.}

\keyformula{\alpha_s(Q^2)=\frac{12\pi}{(33-2n_f)\ln(Q^2/\Lambda^2)},\qquad V(r)=-\frac43\frac{\alpha_s\hbar c}{r}+\kappa r}
 {$n_f$ — число активных ароматов, $\Lambda\approx0{,}2$--$0{,}3$~ГэВ, $\kappa\approx1$~ГэВ/фм.}
 {первая формула — однопетлевое приближение при $Q\gg\Lambda$; вторая — феноменологическая.}
 {асимптотическая свобода и конфайнмент; спектры кваркония.}

\section{Физическая интерпретация}

\textbf{Два режима одного взаимодействия.} На малых расстояниях ($\ll1$~фм) кварки ведут себя почти как свободные
(асимптотическая свобода — партонная картина глубоконеупругого рассеяния), на больших — притяжение растёт
(конфайнмент). Поэтому кварки «видны» в столкновениях высоких энергий, но никогда не вылетают поодиночке.

\textbf{Сравнение с КЭД.} Структура теорий подобна (калибровочные бозоны со спином 1), но цветной заряд глюонов
меняет знак «бега» константы. В КЭД нет конфайнмента: электрон и позитрон можно разделить.

\textbf{Иерархия масштабов.} Сильные распады ($\Delta\to N\pi$) идут за $\sim10^{-23}$~с ($\Gamma\approx120$~МэВ). Слабые распады
странных частиц ($\Lambda\to p\pi^-$) — за $\sim10^{-10}$~с, потому что сильное взаимодействие сохраняет странность, а
изменить аромат может лишь слабое (\tref{t:16}).

\section{Что необходимо знать на экзамене}

\begin{examcheck}
\checkhead{Определения}
\begin{itemize}
 \item адроны, мезоны, барионы; кварковый состав $p,n,\pi^\pm,\pi^0,K^+,\Lambda,\Delta^{++},\Omega^-$;
 \item цвет, цветовой синглет, конфайнмент, асимптотическая свобода, глюоны (8, цветные).
\end{itemize}
\checkhead{Формулы}
\begin{itemize}
 \item заряды и барионные числа кварков; $R=N_c\sum Q_q^2$;
 \item $V(r)=-\frac43\alpha_s\hbar c/r+\kappa r$; качественный вид $\alpha_s(Q^2)$ (убывает логарифмически).
\end{itemize}
\checkhead{Выводы}
\begin{itemize}
 \item проверка зарядов и спинов адронов; проблема $\Delta^{++}$ и её решение цветом;
 \item вычисление $R$ для $u,d,s$ / $+c$ / $+b$ и вывод $N_c=3$.
\end{itemize}
\checkhead{Интерпретации}
\begin{itemize}
 \item эксперименты: глубоконеупругое рассеяние, струи, $J/\psi$; предсказание $\Omega^-$;
 \item почему одиночные кварки не наблюдаются (разрыв струны); откуда масса протона; ядерные силы как остаточные.
\end{itemize}
\end{examcheck}
